\IfFileExists{stacks-project.cls}{%
\documentclass{stacks-project}
}{%
\documentclass{amsart}
}
\usepackage{amssymb}

% For dealing with references we use the comment environment
\usepackage{verbatim}
\usepackage{fancyvrb}
\newenvironment{reference}{\comment}{\endcomment}
%\newenvironment{reference}{}{}
\newenvironment{slogan}{\comment}{\endcomment}
\newenvironment{history}{\comment}{\endcomment}

% For source figures retained as assets
\usepackage{graphicx}

% For commutative diagrams we use Xy-pic
\usepackage[all]{xy}

% We use 2cell for 2-commutative diagrams.
\xyoption{2cell}
\UseAllTwocells

% We use multicol for the list of chapters between chapters
\usepackage{multicol}

% This is generally recommended for better output
\usepackage{lmodern}
\usepackage[T1]{fontenc}
\usepackage{textalpha}

% For cross-file-references
% We prefer the snapshot tmp/xref/#2.aux,
% which latexmkrc replaces atomically, so
% that documents compiled simultaneously
% never read a half-written .aux file.
% The last argument keeps links to #2.pdf.
\usepackage{xr-hyper}
\newcommand{\maybeexternaldocument}[2]{%
  \IfFileExists{#2.aux}{%
    \externaldocument[#1]{#2}%
  }{%
    \PackageWarning{stack}{Missing #2.aux; external references with prefix #1 unavailable}%
  }%
}
\let\externaldocumentorig\externaldocument
\renewcommand{\externaldocument}[2][]{%
  \IfFileExists{tmp/xref/#2.aux}{%
    \externaldocumentorig[#1]{tmp/xref/#2}[#2.pdf]%
  }{%
  \IfFileExists{#2.aux}{%
    \externaldocumentorig[#1]{#2}%
  }{%
    \PackageWarning{stack}{Missing #2.aux; external references with prefix #1 unavailable}%
  }%
  }%
}

% Package for hypertext links:
\usepackage{hyperref}

% For any local file, say "hello.tex" you want to link to please
% use \externaldocument[hello-]{hello}
\externaldocument[introduction-]{introduction}
\externaldocument[conventions-]{conventions}
\externaldocument[sets-]{sets}
\externaldocument[categories-]{categories}
\externaldocument[topology-]{topology}
\externaldocument[sheaves-]{sheaves}
\externaldocument[sites-]{sites}
\externaldocument[stacks-]{stacks}
\externaldocument[fields-]{fields}
\externaldocument[algebra-]{algebra}
\externaldocument[brauer-]{brauer}
\externaldocument[homology-]{homology}
\externaldocument[derived-]{derived}
\externaldocument[simplicial-]{simplicial}
\externaldocument[more-algebra-]{more-algebra}
\externaldocument[smoothing-]{smoothing}
\externaldocument[modules-]{modules}
\externaldocument[sites-modules-]{sites-modules}
\externaldocument[injectives-]{injectives}
\externaldocument[cohomology-]{cohomology}
\externaldocument[sites-cohomology-]{sites-cohomology}
\externaldocument[dga-]{dga}
\externaldocument[dpa-]{dpa}
\externaldocument[sdga-]{sdga}
\externaldocument[hypercovering-]{hypercovering}
\externaldocument[schemes-]{schemes}
\externaldocument[constructions-]{constructions}
\externaldocument[properties-]{properties}
\externaldocument[morphisms-]{morphisms}
\externaldocument[coherent-]{coherent}
\externaldocument[divisors-]{divisors}
\externaldocument[limits-]{limits}
\externaldocument[varieties-]{varieties}
\externaldocument[topologies-]{topologies}
\externaldocument[descent-]{descent}
\externaldocument[perfect-]{perfect}
\externaldocument[more-morphisms-]{more-morphisms}
\externaldocument[flat-]{flat}
\externaldocument[groupoids-]{groupoids}
\externaldocument[more-groupoids-]{more-groupoids}
\externaldocument[etale-]{etale}
\externaldocument[chow-]{chow}
\externaldocument[intersection-]{intersection}
\externaldocument[pic-]{pic}
\externaldocument[weil-]{weil}
\externaldocument[adequate-]{adequate}
\externaldocument[dualizing-]{dualizing}
\externaldocument[duality-]{duality}
\externaldocument[discriminant-]{discriminant}
\externaldocument[derham-]{derham}
\externaldocument[local-cohomology-]{local-cohomology}
\externaldocument[algebraization-]{algebraization}
\externaldocument[curves-]{curves}
\externaldocument[resolve-]{resolve}
\externaldocument[models-]{models}
\externaldocument[functors-]{functors}
\externaldocument[equiv-]{equiv}
\externaldocument[pione-]{pione}
\externaldocument[etale-cohomology-]{etale-cohomology}
\externaldocument[proetale-]{proetale}
\externaldocument[relative-cycles-]{relative-cycles}
\externaldocument[more-etale-]{more-etale}
\externaldocument[trace-]{trace}
\externaldocument[crystalline-]{crystalline}
\externaldocument[spaces-]{spaces}
\externaldocument[spaces-properties-]{spaces-properties}
\externaldocument[spaces-morphisms-]{spaces-morphisms}
\externaldocument[decent-spaces-]{decent-spaces}
\externaldocument[spaces-cohomology-]{spaces-cohomology}
\externaldocument[spaces-limits-]{spaces-limits}
\externaldocument[spaces-divisors-]{spaces-divisors}
\externaldocument[spaces-over-fields-]{spaces-over-fields}
\externaldocument[spaces-topologies-]{spaces-topologies}
\externaldocument[spaces-descent-]{spaces-descent}
\externaldocument[spaces-perfect-]{spaces-perfect}
\externaldocument[spaces-more-morphisms-]{spaces-more-morphisms}
\externaldocument[spaces-flat-]{spaces-flat}
\externaldocument[spaces-groupoids-]{spaces-groupoids}
\externaldocument[spaces-more-groupoids-]{spaces-more-groupoids}
\externaldocument[bootstrap-]{bootstrap}
\externaldocument[spaces-pushouts-]{spaces-pushouts}
\externaldocument[spaces-chow-]{spaces-chow}
\externaldocument[groupoids-quotients-]{groupoids-quotients}
\externaldocument[spaces-more-cohomology-]{spaces-more-cohomology}
\externaldocument[spaces-simplicial-]{spaces-simplicial}
\externaldocument[spaces-duality-]{spaces-duality}
\externaldocument[formal-spaces-]{formal-spaces}
\externaldocument[restricted-]{restricted}
\externaldocument[spaces-resolve-]{spaces-resolve}
\externaldocument[formal-defos-]{formal-defos}
\externaldocument[defos-]{defos}
\externaldocument[cotangent-]{cotangent}
\externaldocument[examples-defos-]{examples-defos}
\externaldocument[algebraic-]{algebraic}
\externaldocument[examples-stacks-]{examples-stacks}
\externaldocument[stacks-sheaves-]{stacks-sheaves}
\externaldocument[criteria-]{criteria}
\externaldocument[artin-]{artin}
\externaldocument[quot-]{quot}
\externaldocument[stacks-properties-]{stacks-properties}
\externaldocument[stacks-morphisms-]{stacks-morphisms}
\externaldocument[stacks-limits-]{stacks-limits}
\externaldocument[stacks-cohomology-]{stacks-cohomology}
\externaldocument[stacks-perfect-]{stacks-perfect}
\externaldocument[stacks-introduction-]{stacks-introduction}
\externaldocument[stacks-more-morphisms-]{stacks-more-morphisms}
\externaldocument[stacks-geometry-]{stacks-geometry}
\externaldocument[moduli-]{moduli}
\externaldocument[moduli-curves-]{moduli-curves}
\externaldocument[examples-]{examples}
\externaldocument[exercises-]{exercises}
\externaldocument[guide-]{guide}
\externaldocument[desirables-]{desirables}
\externaldocument[coding-]{coding}
\externaldocument[obsolete-]{obsolete}
\externaldocument[fdl-]{fdl}
\externaldocument[index-]{index}

% Heap Project documents
\externaldocument[k3-]{k3}
\externaldocument[torelli-]{torelli}
\externaldocument[linear-algebra-]{linear-algebra}
\externaldocument[groups-]{groups}
\externaldocument[complex-analysis-]{complex-analysis}
\externaldocument[functional-analysis-]{functional-analysis}
\externaldocument[categories-]{categories}
\externaldocument[commutative-algebra-]{commutative-algebra}
\externaldocument[ringed-spaces-]{ringed-spaces}
\externaldocument[homological-algebra-]{homological-algebra}
\externaldocument[lie-algebras-]{lie-algebras}
\externaldocument[representation-theory-]{representation-theory}
\externaldocument[smooth-manifolds-]{smooth-manifolds}
\externaldocument[differential-topology-]{differential-topology}
\externaldocument[algebraic-topology-]{algebraic-topology}
\externaldocument[differential-geometry-]{differential-geometry}
\externaldocument[algebraic-geometry-]{algebraic-geometry}
\externaldocument[symplectic-geometry-]{symplectic-geometry}
\externaldocument[complex-geometry-]{complex-geometry}
\externaldocument[hodge-theory-]{hodge-theory}
\externaldocument[%
variations-of-hodge-structures-]{%
variations-of-hodge-structures}
\externaldocument[%
homological-mirror-symmetry-]{%
homological-mirror-symmetry}
\maybeexternaldocument{deformations-}{deformations}
\externaldocument[vertex-algebras-]{vertex-algebras}
\externaldocument[springer-]{springer}
\externaldocument[infty-categories-]{infty-categories}
\externaldocument[seminars-differential-geometry-]{%
seminars-differential-geometry}
\externaldocument[seminars-complex-geometry-]{%
seminars-complex-geometry}
\externaldocument[seminars-algebraic-geometry-]{%
seminars-algebraic-geometry}
\externaldocument[seminars-others-]{%
seminars-others}
\externaldocument[%
geometric-prequantization-]{%
geometric-prequantization}
\externaldocument[surfaces-]{surfaces}
\externaldocument[birational-maps-commutative-algebra-]{birational-maps-commutative-algebra}
\externaldocument[cremona-transformations-]{cremona-transformations}
\externaldocument[derived-categories-of-sheaves-]{derived-categories-of-sheaves}
\externaldocument[geometric-invariant-theory-]{geometric-invariant-theory}
\externaldocument[geometric-stability-conditions-and-group-actions-]{geometric-stability-conditions-and-group-actions}
\externaldocument[moduli-spaces-of-sheaves-]{moduli-spaces-of-sheaves}
\externaldocument[bridgeland-stability-general-theory-]{bridgeland-stability-general-theory}
\externaldocument[hilbert-schemes-]{hilbert-schemes}
\externaldocument[noncommutative-abelian-surfaces-and-kummer-type-hyperkahler-varieties-]{noncommutative-abelian-surfaces-and-kummer-type-hyperkahler-varieties}
\externaldocument[mumford-tate-groups-in-hodge-theory-]{mumford-tate-groups-in-hodge-theory}
\externaldocument[hodge-birational-atoms-2026-]{hodge-birational-atoms-2026}
\externaldocument[xxii-egd-2026-]{%
xxii-egd-2026}
\externaldocument[pde-]{pde}
\externaldocument[probability-]{probability}

\externaldocument[linguistics-]{linguistics}
\externaldocument[genealogia-familiar-]{%
genealogia-familiar}

\externaldocument[%
16th-alga-meeting-2026-]{%
16th-alga-meeting-2026}

\externaldocument[%
iv-jornadas-lefschetz-2026-]{%
iv-jornadas-lefschetz-2026}

\externaldocument[reading-the-masters-]{%
reading-the-masters}

\externaldocument[real-analysis-]{%
real-analysis}

\externaldocument[optimal-transport-]{%
optimal-transport}

\externaldocument[history-]{history}

\externaldocument[seiberg-witten-]{seiberg-witten}

\externaldocument[history-philosophy-biology-]{history-philosophy-biology}

\externaldocument[chemistry-]{chemistry}

\externaldocument[molecular-cell-biology-]{molecular-cell-biology}

\externaldocument[microbiology-]{microbiology}

\externaldocument[protistology-]{protistology}

\externaldocument[mycology-]{mycology}

\externaldocument[botany-]{botany}

\externaldocument[zoology-]{zoology}

\externaldocument[animal-histology-]{animal-histology}

\externaldocument[animal-physiology-]{animal-physiology}

\externaldocument[systematics-biogeography-]{systematics-biogeography}

\externaldocument[ecology-natural-resources-]{ecology-natural-resources}

\externaldocument[earth-sciences-]{earth-sciences}

\externaldocument[biostatistics-]{biostatistics}

\externaldocument[mathematical-systems-biology-]{mathematical-systems-biology}

\externaldocument[cimpa-el-salvador-2024-]{cimpa-el-salvador-2024}

% Theorem environments.
%
\theoremstyle{plain}
\newtheorem{theorem}[subsection]{Theorem}
\newtheorem{proposition}[subsection]{Proposition}
\newtheorem{lemma}[subsection]{Lemma}
\newtheorem{corollary}[subsection]{Corollary}

\theoremstyle{definition}
\newtheorem{definition}[subsection]{Definition}
\newtheorem{example}[subsection]{Example}
\newtheorem{exercise}[subsection]{Exercise}
\newtheorem{situation}[subsection]{Situation}

\theoremstyle{remark}
\newtheorem{remark}[subsection]{Remark}
\newtheorem{remarks}[subsection]{Remarks}

\numberwithin{equation}{subsection}

% Macros
%
\def\lim{\mathop{\mathrm{lim}}\nolimits}
\def\colim{\mathop{\mathrm{colim}}\nolimits}
\def\Spec{\mathop{\mathrm{Spec}}}
\def\Hom{\mathop{\mathrm{Hom}}\nolimits}
\def\Ext{\mathop{\mathrm{Ext}}\nolimits}
\def\SheafHom{\mathop{\mathcal{H}\!\mathit{om}}\nolimits}
\def\SheafExt{\mathop{\mathcal{E}\!\mathit{xt}}\nolimits}
\def\Sch{\mathit{Sch}}
\def\Mor{\mathop{\mathrm{Mor}}\nolimits}
\def\Ob{\mathop{\mathrm{Ob}}\nolimits}
\def\Sh{\mathop{\mathit{Sh}}\nolimits}
\def\NL{\mathop{N\!L}\nolimits}
\def\CH{\mathop{\mathrm{CH}}\nolimits}
\def\proetale{{pro\text{-}\acute{e}tale}}
\def\etale{{\acute{e}tale}}
\def\QCoh{\mathit{QCoh}}
\def\Ker{\mathop{\mathrm{Ker}}}
\def\Im{\mathop{\mathrm{Im}}}
\def\Coker{\mathop{\mathrm{Coker}}}
\def\Coim{\mathop{\mathrm{Coim}}}

% Boxtimes
%
\DeclareMathSymbol{\boxtimes}{\mathbin}{AMSa}{"02}

%
% Macros for moduli stacks/spaces
%
\def\QCohstack{\mathcal{QC}\!\mathit{oh}}
\def\Cohstack{\mathcal{C}\!\mathit{oh}}
\def\Spacesstack{\mathcal{S}\!\mathit{paces}}
\def\Quotfunctor{\mathrm{Quot}}
\def\Hilbfunctor{\mathrm{Hilb}}
\def\Curvesstack{\mathcal{C}\!\mathit{urves}}
\def\Polarizedstack{\mathcal{P}\!\mathit{olarized}}
\def\Complexesstack{\mathcal{C}\!\mathit{omplexes}}
% \Pic is the operator that assigns to X its picard group, usage \Pic(X)
% \Picardstack_{X/B} denotes the Picard stack of X over B
% \Picardfunctor_{X/B} denotes the Picard functor of X over B
\def\Pic{\mathop{\mathrm{Pic}}\nolimits}
\def\Picardstack{\mathcal{P}\!\mathit{ic}}
\def\Picardfunctor{\mathrm{Pic}}
\def\Deformationcategory{\mathcal{D}\!\mathit{ef}}


\begin{document}

\title{Earth Sciences}
\maketitle

\phantomsection
\label{section-phantom}

\tableofcontents

\section{What this chapter covers}
\label{section-what-this-chapter-covers}

Brief notes on physical geology: the
structure of the Earth, plate tectonics,
rocks and the processes that make and
destroy them, geologic time, the ocean
floor, the agents that shape the surface,
resources and the other planets. Links to
biology are pointed out where they arise:
fossils and dating, biogeography and
moving continents, soils, rocks made by
organisms, and the effect of glaciations on
living things.

\textbf{Sources of these notes.}
The course Ciencias de la Tierra (second
semester, spring 2014) at the Facultad de
Ciencias, UNAM. Its textbook was Tarbuck and
Lutgens, \emph{Ciencias de la Tierra}
\cite{tarbuckciencias}, chapters 1--22.
What survives is Daniel's typed answers to
the questionnaires on chapters 1--13 and
15--22, and scans of chapters 1--3 of the
textbook. The notes follow his answers,
corrected and completed with the scanned
chapters and standard facts.

\section{The Earth system and the rock
cycle}
\label{section-earth-system}

\textbf{Geology.} \emph{Physical geology}
studies Earth materials and the processes
acting on and below the surface;
\emph{historical geology} studies the
origin of the Earth and its changes through
time, including the history of life. Against
catastrophism (and Ussher's creation date of
4004 BC), Hutton's \emph{uniformitarianism}
holds that the processes acting today also
acted in the past, so that the present is
the key to the past. The Earth is about
$4.5$ billion years old
\cite[Chapter 1]{tarbuckciencias}.

\textbf{Spheres.} The Earth is treated as a
system of four interacting parts:
hydrosphere, atmosphere, biosphere and
geosphere (solid Earth). It is powered by
the Sun (external processes such as
weather, erosion) and by internal heat
(volcanism, mountain building). As in
physiology, \emph{negative feedback} damps a
change and \emph{positive feedback}
amplifies it.

\textbf{Origin.} In the nebular hypothesis a
cloud of gas and dust contracted about
$5$ billion years ago; most matter
collected into the protosun, and the
planets accreted from the remaining disk,
the rocky inner planets close to the Sun
and the gas- and ice-rich outer planets far
from it.

\textbf{Layers.} By composition: crust
(oceanic about $7$ km thick, basaltic;
continental $35$--$40$ km on average,
granitic, density about
$2.7\ \mathrm{g/cm^3}$), mantle (to
$2{,}900$ km depth, mostly peridotite, more
than $82\%$ of the Earth's volume) and core
(iron--nickel, radius $3{,}486$ km). By
mechanical behaviour: rigid
\emph{lithosphere} (crust and uppermost
mantle, $5$--$250$ km), weak
\emph{asthenosphere} (to $660$ km), on
which the lithosphere moves, mesosphere,
liquid outer core and solid inner core
(radius about $1{,}216$ km)
\cite[Chapter 1]{tarbuckciencias}.

\textbf{The rock cycle.} The three rock
classes turn into each other:
\begin{itemize}
\item \emph{igneous} rocks crystallize from
magma, either inside the crust
(\emph{plutonic}, e.g. granite) or at the
surface (\emph{volcanic}, e.g. basalt);
\item weathering and erosion break rocks
into sediment, which is compacted and
cemented (lithified) into
\emph{sedimentary} rock, e.g. sandstone,
shale;
\item heat, pressure and fluids turn any
rock into \emph{metamorphic} rock, e.g.
slate, gneiss, marble (from limestone);
\item metamorphic rock that melts gives
magma again.
\end{itemize}
Every rock is made from the material of an
older rock; shortcuts are possible (an
igneous rock can be weathered or
metamorphosed directly).

\section{Plate tectonics}
\label{section-plate-tectonics}

\textbf{Continental drift.} Wegener proposed
that the continents once formed one
supercontinent, \emph{Pangaea}, which began
to break up about $200$ million years ago.
His evidence
\cite[Chapter 2]{tarbuckciencias}:
\begin{itemize}
\item the fit of South America and Africa,
best along the continental slope at about
$900$ m depth;
\item fossils: \emph{Mesosaurus}, an
aquatic reptile of the Permian (about
$260$ million years ago), is found only in
eastern South America and southern Africa;
the seed fern \emph{Glossopteris}, with
large seeds hard to disperse, occurs in the
late Paleozoic of Africa, Australia, India,
South America and Antarctica;
\item matching rock types and mountain
belts across the Atlantic (the
Appalachians and the Caledonian mountains
of the British Isles and Scandinavia);
\item palaeoclimate: late Paleozoic glacial
deposits in now tropical southern
continents, explained if they lay together
near the South Pole.
\end{itemize}
The alternative of the time, sunken land
bridges between continents, is ruled out:
no trace of them exists on the ocean floor.

\textbf{Palaeomagnetism and sea-floor
spreading.} Iron-rich minerals such as
magnetite record the Earth's field when a
rock cools; the inclination of that record
gives the latitude at which the rock
formed. Hess proposed \emph{sea-floor
spreading}: mantle material rises at
mid-ocean ridges, spreads sideways and cools
into new ocean floor. Vine and Matthews
connected this with the stripes of normal
and reversed magnetization parallel to the
ridges, symmetric on both sides: each
reversal of the Earth's field is recorded
as the new floor moves away. The oldest
continental rocks are about $4$ billion
years old, the oldest sea floor less than
$200$ million.

\textbf{Plate boundaries.} The lithosphere
is divided into plates.
\begin{itemize}
\item \emph{Divergent}: plates separate and
new lithosphere forms (mid-ocean ridges,
continental rifts such as East Africa).
\item \emph{Convergent}: lithosphere is
consumed. The denser oceanic plate
subducts under continental or older
oceanic lithosphere, giving trenches,
volcanic arcs and deep earthquakes. When
two continents meet, neither subducts far
and they collide: the Himalaya rose when
India met Asia, about $45$ million years
ago \cite[Chapter 2]{tarbuckciencias}.
\item \emph{Transform}: plates slide past
each other along a fault, neither creating
nor destroying lithosphere.
\end{itemize}
Since the Earth does not grow, lithosphere
created at ridges must be consumed at
subduction zones.

\begin{remark}
\label{remark-tectonics-biogeography}
Plate tectonics is the physical basis of
historical biogeography. Disjunct
distributions like those of
\emph{Mesosaurus} and \emph{Glossopteris}
record former connections, and the
break-up of Pangaea split biotas apart
(vicariance). Wegener also cited living
organisms, notably the Australian
marsupials, whose isolated evolution
follows the long separation of Australia
\cite[Chapter 2]{tarbuckciencias}.
\end{remark}

\section{Minerals, igneous rocks and
volcanoes}
\label{section-igneous}

\begin{definition}
\label{definition-mineral}
A \emph{mineral} is a naturally occurring,
inorganic solid with an ordered internal
structure and a definite chemical
composition (within limits). A solid
lacking the ordered structure, such as
opal, is a \emph{mineraloid}. A
\emph{rock} is a natural solid mass of
minerals or mineral-like matter.
\end{definition}

Atoms of the same element with different
numbers of neutrons are \emph{isotopes}.
Ionic bonds come from transfer of
electrons, covalent bonds from sharing of
electron pairs. Colour is a poor diagnostic
of minerals (impurities change it);
hardness, cleavage, streak and crystal form
are better.

\textbf{Igneous rocks} are classified by
texture and composition. Texture records
the cooling history:
\begin{itemize}
\item \emph{aphanitic} (fine-grained, fast
cooling at the surface), \emph{phaneritic}
(coarse-grained, slow cooling at depth),
\emph{porphyritic} (large crystals in a
fine matrix: two cooling stages),
\emph{glassy} (quenched, e.g. obsidian),
\emph{pegmatitic} (crystals over $1$ cm,
grown from fluid-rich magma);
\item composition ranges from
\emph{felsic} (granitic: rich in silica,
quartz and feldspar, light) through
intermediate (andesitic) and \emph{mafic}
(basaltic) to ultramafic (peridotite).
Granite and rhyolite have the same
composition, one plutonic, one volcanic.
\end{itemize}

\textbf{Magma.} Rocks melt when temperature
rises (the geothermal gradient), when
pressure drops (\emph{decompression
melting} under ridges), or when water is
added (subduction zones). In \emph{Bowen's
reaction series} the mafic minerals
crystallize first, at high temperature
(olivine, pyroxene, amphibole, biotite,
along with calcium-rich plagioclase), and
the felsic ones last (potassium feldspar,
muscovite, quartz). Removing early crystals
changes the remaining melt
(\emph{magmatic differentiation}), and
\emph{partial melting} of a rock yields a
melt richer in silica than its source.
Most basaltic magma comes from partial
melting of mantle peridotite; andesitic and
granitic magmas need differentiation or
melting of crust.

\textbf{Volcanoes.} Explosiveness is
governed by viscosity, which grows with
silica content and falls with temperature,
and by dissolved gas, mostly water vapour
(about $70\%$), then CO2 (about $15\%$).
Viscous, gas-rich magma traps gas and
erupts explosively; fluid basaltic lava
lets gas escape quietly.
\begin{itemize}
\item \emph{Shield volcanoes}: broad, of
fluid basalt, quiet eruptions (Kilauea).
\item \emph{Cinder cones}: small, steep,
built of ejected fragments (Paricut\'in,
Michoac\'an).
\item \emph{Composite cones}: lava and
pyroclastic layers, andesitic, explosive
(Mount St.\ Helens, 1980).
\end{itemize}
Plutonic bodies include \emph{dikes}
(cutting across layers), \emph{sills}
(parallel to them), \emph{laccoliths}
(arching the layers above) and
\emph{batholiths} (the largest). Volcanic
outgassing produced much of the early
atmosphere and oceans.

\section{Weathering, soils, sedimentary and
metamorphic rocks}
\label{section-surface-rocks}

\textbf{Weathering.} \emph{Mechanical}
weathering breaks rock without changing its
composition (frost wedging, unloading and
exfoliation domes); \emph{chemical}
weathering changes it. Breaking rock
increases the surface on which chemical
weathering acts. Carbonic acid, formed from
CO2 dissolved in rain or in soil water,
turns potassium feldspar into kaolinite:
$$
\mathrm{2KAlSi_3O_8 + 2H^+ + H_2O
\rightarrow Al_2Si_2O_5(OH)_4
+ 2K^+ + 4SiO_2}.
$$
Chemical weathering is slow in dry climates
for lack of water.

\textbf{Soil} is the part of the regolith
that supports plants: mineral matter,
organic matter, water and air. Its
formation depends on parent rock, climate,
topography (slope and its orientation),
organisms and time. The \emph{horizons}:
\begin{itemize}
\item O: organic matter, litter and living
organisms;
\item A: mineral matter mixed with humus,
high biological activity;
\item E: pale, leached zone (eluviation);
\item B: accumulation of material washed
down from above;
\item C: partly weathered parent rock.
\end{itemize}

\textbf{Sedimentary rocks} form a thin skin
(about $5\%$ of the outer $16$ km of
crust) but about $75\%$ of continental
outcrops. \emph{Detrital} rocks are
classified by grain size:
conglomerate (rounded) and breccia
(angular), sandstone, siltstone, shale.
Quartz and clay dominate them, quartz
because it resists weathering, clay because
it is the product of weathering of
feldspars. \emph{Chemical} rocks
precipitate from solution, inorganically
(evaporites such as halite, travertine) or
\emph{biochemically}: limestone from shells
of foraminifera and corals, chert, and coal
from plant matter. The hallmark of
sedimentary rocks is \emph{bedding}, with
cross-bedding, graded beds, ripple marks
(asymmetric for currents, symmetric for
waves), and fossils.

\textbf{Metamorphism} is driven by heat (the
main agent), pressure (confining or
\emph{differential}) and chemically active
fluids, acting on a parent rock
(protolith). Differential stress aligns
platy minerals into \emph{foliation}. With
increasing grade, shale becomes
\begin{itemize}
\item \emph{slate} (very fine, splits into
flat sheets),
\item \emph{phyllite} (satiny sheen, wavy
surfaces),
\item \emph{schist} (visible micas),
\item \emph{gneiss} (banded),
\end{itemize}
and partial melting of gneiss gives
migmatite. Non-foliated rocks include
marble and quartzite. Settings: contact
(around intrusions), regional (convergent
boundaries, the cores of mountain belts),
hydrothermal, burial, and along faults.
\emph{Index minerals} indicate the grade.

\section{Geologic time and dating}
\label{section-geologic-time}

\textbf{Relative dating} orders events
without assigning years:
\begin{itemize}
\item \emph{superposition}: in undeformed
strata each layer is younger than the one
below and older than the one above;
\item \emph{original horizontality}:
sediments are deposited in horizontal
layers;
\item \emph{cross-cutting}: a dike or fault
is younger than the rocks it cuts;
\item \emph{inclusions}: fragments enclosed
in a rock are older than the rock;
\item \emph{unconformities} record gaps:
angular unconformity (tilted beds below,
horizontal above), disconformity (parallel
beds), nonconformity (sediments on igneous
or metamorphic rock).
\end{itemize}
Rocks of the same age in different places
are matched by \emph{correlation}.

\textbf{Fossils.} William Smith found that
each formation contains a characteristic
set of fossils, distinct from those above
and below: \emph{faunal succession}.
Fossilization favours organisms with hard
parts and rapid burial. Fossils include
molds and casts, carbonized remains,
replacement, organisms in amber, and trace
fossils (tracks, burrows, coprolites,
gastroliths). Widespread species of short
duration are the best index fossils.

\textbf{Numerical dating.} A radioactive
isotope decays with constant half-life
$T$: after time $t$ a fraction
$$
\frac{N(t)}{N_0} = 2^{-t/T},
\qquad
t = T\log_2\frac{N_0}{N(t)}
$$
remains, provided the system stayed closed
(no gain or loss of parent or daughter).
Tree rings give annual resolution for
recent times.

\textbf{The time scale} is divided into
eons, eras, periods, epochs and ages. The
Precambrian is poorly subdivided: fossils
with hard parts appear only at the start
of the Cambrian, and old rocks have been
deformed and metamorphosed.

\section{Deformation, earthquakes and the
interior}
\label{section-earthquakes-interior}

\textbf{Deformation.} \emph{Stress} is force
per unit area (compressional, tensional,
shear); \emph{strain} is the resulting
change of shape or volume. Beyond the
elastic limit rocks deform
\emph{ductilely} (fold or flow), favoured by
high temperature, high confining pressure
and slow loading, or \emph{brittlely}
(fracture) near the surface. Orientation of
a planar layer is given by its
\emph{strike} (direction of its
intersection with a horizontal plane) and
\emph{dip} (inclination).
\begin{itemize}
\item Folds: \emph{anticlines} (arched up),
\emph{synclines} (down); domes and basins.
\item Faults: normal (tension; horsts and
grabens, as in rift valleys), reverse and
low-angle \emph{thrust} faults
(compression), strike-slip; joints are
fractures without displacement.
\end{itemize}

\textbf{Earthquakes} are vibrations caused
by sudden slip on a fault, released from
the \emph{focus} (the epicentre is above it).
Elastic rebound: rocks bend under stress,
then slip and spring back. Waves:
\begin{itemize}
\item \emph{P waves}, compressional, the
fastest, travel through solids and
liquids;
\item \emph{S waves}, shear, slower, travel
only through solids;
\item surface waves, the most damaging.
\end{itemize}
Most large earthquakes occur at subduction
zones; deep-focus earthquakes occur below
the trenches. \emph{Intensity} (Mercalli)
measures effects at a place;
\emph{magnitude} measures the source. On
the Richter scale one unit means $10$ times
the wave amplitude and about $32$ times the
energy; the moment magnitude, computed from
fault area and slip, is better for large
earthquakes. Damage depends also on ground
and building design: in the 1985 Mexico
City earthquake the soft lake-bed sediments
under the city amplified the shaking.

\textbf{The interior.} Seismic velocities
depend on density and elasticity, so
reflections and refractions map the
layers. Mohorovi\v{c}i\'c found a jump in
velocity at the base of the crust (the
\emph{Moho}), a change of composition. The
outer core does not transmit S waves, so it
is liquid. Internal heat comes from decay
of uranium, thorium and potassium, from
the Earth's formation, and from the
crystallization of the inner core.

\section{The ocean floor}
\label{section-ocean-floor}

Depth is measured by echo sounding: with
sound speed $1{,}500$ m/s and a return time
of $6$ s, the depth is
$\frac12(1{,}500\cdot 6) = 4{,}500$ m.
Satellites measure the sea surface by radar
altimetry; it mirrors the relief of the
floor.

\textbf{Provinces.}
\begin{itemize}
\item \emph{Continental margins}: shelf,
slope and rise. \emph{Passive} margins (most
of the Atlantic) lie within a plate;
\emph{active} margins (around the Pacific)
lie over subduction zones, with an
accretionary wedge of scraped-off sediment.
\item \emph{Deep-ocean basins}: abyssal
plains (larger in the Atlantic, which has
few trenches to trap sediment), trenches,
seamounts and flat-topped guyots.
\item \emph{Mid-ocean ridges}: elevated
because the young lithosphere is hot and
less dense; it cools, thickens and subsides
with age.
\end{itemize}

Oceanic crust, seen in ophiolites, has
layers: deep-sea sediment, pillow basalts,
sheeted dikes, gabbro. At ridges seawater
circulates through the hot crust, alters
the basalt and returns laden with dissolved
metals, precipitating at hydrothermal vents
(black smokers), which host communities
living on chemosynthesis.

Continental rifting (East Africa) may open
a new ocean basin; mantle plumes may help
start rifts but are too narrow to drive
plates. Cold, dense oceanic lithosphere
sinks at subduction zones. The Farallon
plate was almost entirely subducted under
the Americas; the Juan de Fuca, Cocos and
Nazca plates are its remnants. Continents
assemble into supercontinents and break up
again (the \emph{supercontinent cycle}).

\section{Surface processes}
\label{section-surface-processes}

\textbf{Mass wasting} moves rock and soil
downslope under gravity. Water reduces
cohesion and adds weight; slopes steeper
than the \emph{angle of repose} are
unstable; removing vegetation (roots) and
earthquakes trigger movements. Types:
\emph{falls}, \emph{slides} along a
surface (slumps on curved surfaces, rock
slides on bedding or joints), \emph{flows}
(mudflows in semi-arid regions, earthflows
in humid ones), and slow \emph{creep}
(tilted fences and trees). In permafrost
regions the thawed surface layer flows
(solifluction).

\textbf{Rivers.} In the hydrologic cycle
water evaporates, precipitates, and on land
infiltrates, runs off, is transpired by
plants or is stored as ice. Ocean
evaporation exceeds precipitation over the
ocean, and the difference returns as
runoff. \emph{Discharge} is width times
depth times velocity. The \emph{gradient}
is the drop per unit length: a stream
falling $2{,}000$ m over $250$ km has mean
gradient $8$ m/km. \emph{Base level} is the
lowest level to which a stream can erode.
\emph{Capacity} (maximum load, set by
discharge) differs from \emph{competence}
(largest grain, set by velocity). Rivers
build natural levees, deltas (where they
enter standing water) and alluvial fans
(where they leave a mountain front). Flood
control by artificial levees, dams and
channelization each has costs.

\textbf{Groundwater} is the water of the
\emph{zone of saturation}; its top is the
water table, which follows the topography
in subdued form. It is most of the liquid
fresh water on land. \emph{Porosity} is the
fraction of open space; \emph{permeability}
is the ability to transmit fluid. Clay is
porous but impermeable, since its pores are
tiny. An \emph{aquifer} transmits water
freely, an \emph{aquitard} or aquiclude
does not. A \emph{gaining} (effluent)
stream is fed by groundwater; a
\emph{losing} (influent) stream lies above
the water table and feeds it.

\emph{Darcy's law} says that the discharge
$Q$ of groundwater through a cross-section
of area $A$ is
\begin{equation}
\label{equation-darcy}
Q = K A \,\frac{h_1-h_2}{L},
\end{equation}
where $(h_1-h_2)/L$ is the
\emph{hydraulic gradient} (drop in
water-table height over the flow distance
$L$) and $K$ the \emph{hydraulic
conductivity}, which measures
permeability. In differential form the
flux per unit area is $q = -K\nabla h$.
Thus flow is faster where the water table
is steeper and the material more
permeable. An \emph{artesian} system needs
an inclined aquifer recharged at one end
and confined above and below by
aquitards. Hot springs and geysers are
mostly heated by cooling igneous rock.
Sandstone filters contaminated water better
than coarse gravel or karstic limestone.

\textbf{Glaciers} cover about $10\%$ of the
land today, about three times more in the
Pleistocene. A glacier advances, retreats
or stands still according to its
\emph{budget}: accumulation minus ablation.
It erodes by plucking and abrasion, turning
V-shaped valleys into U-shaped ones with
hanging valleys, cirques, tarns, ar\^etes
and horns, and it deposits till (directly)
and stratified drift (by meltwater):
moraines (lateral, medial, end, ground),
kettles, kames and eskers. Medial moraines
prove that valley glaciers move. The
Pleistocene glaciations lowered sea level,
opening land bridges, and drove migrations
and extinctions of plants and animals.

\textbf{Deserts.} Dry regions cover about
$30\%$ of the land. Subtropical deserts lie
under the subtropical high-pressure belts,
where air subsides and warms; mid-latitude
deserts lie far from the ocean or in rain
shadows. Running water is still the main
erosive agent, though drainage is mostly
interior. Wind moves sand by
\emph{saltation}, rarely more than about
$1$ m high; deflation hollows are limited
by the water table; dunes migrate by
saltation up the gentle side and slipping
down the steep face; \emph{loess} is
wind-blown silt from deserts and glacial
outwash.

\textbf{Coasts.} Wave size depends on wind
speed, duration and \emph{fetch}. In
shallow water waves slow, shorten, steepen
and break. Waves erode by impact and
abrasion and move sand along the shore
(beach drift, longshore currents). Features:
wave-cut cliffs and platforms, uplifted
marine terraces, spits, bay barriers,
tombolos, barrier islands. Groins, jetties
and breakwaters trap sand on one side and
starve the other; beach nourishment and
relocation are the alternatives.

\begin{remark}
\label{remark-tides}
Tides are raised by the Moon and, less, by
the Sun. The tidal force is the difference
of gravitational pull across the Earth, so
it scales as $M/d^3$, not $M/d^2$. With
$M_\odot/M_{\mathrm{Moon}}\approx
2.7\cdot 10^{7}$ and
$d_\odot/d_{\mathrm{Moon}}\approx 389$ this
gives a solar effect about $0.46$ of the
lunar one. Tidal friction slows the
Earth's rotation; growth bands of fossil
corals and bivalves record more days per
year in the past.
\end{remark}

\section{Energy and mineral resources}
\label{section-resources}

\emph{Renewable} resources regenerate in
months to decades (food, wood, solar,
wind); \emph{non-renewable} ones, such as
coal, oil and natural gas, take millions
of years. Oil and gas accumulate in
\emph{traps}: a porous, permeable reservoir
rock under an impermeable cap rock. Burning
sulfur-bearing coal produces sulfur oxides
and acid rain. Nuclear fission uses
uranium-235; dams silt up; geothermal
fields can be exhausted locally.

A \emph{resource} is a deposit, identified
or not, that may be extracted now or later;
a \emph{reserve} is an identified deposit
that is economically extractable now. An
\emph{ore} is a deposit worth mining; a
deposit becomes ore when prices rise or
technology improves. Ore-forming processes:
\begin{itemize}
\item hydrothermal deposits, as veins or
disseminated;
\item weathering: bauxite, the ore of
aluminium, forms by intense tropical
weathering that leaches everything but
aluminium; leaching also concentrates
copper at depth below a rusty iron-oxide
cap;
\item placers, concentrated by running
water: gold, platinum, diamonds, tin.
\end{itemize}
Non-metallic resources (building materials
such as sand, gravel and cement, and
industrial minerals) exceed metals in
value.

\section{Planetary geology}
\label{section-planetary-geology}

The \emph{terrestrial} planets are small,
dense and rocky; the \emph{Jovian} planets
are large, of low density, made of gases
and ices, and rotate fast. The three kinds
of planetary material (rocks, ices, gases)
differ in melting point, which explains the
pattern through condensation in the solar
nebula. Atmospheres depend on mass (escape
velocity), temperature and, on Earth, life,
which turned the atmosphere oxygen-rich.

\begin{itemize}
\item \emph{Moon}: surfaces with more
craters are older (crater density is a
relative clock); the dark \emph{maria} are
floods of fluid basalt, like the Columbia
Plateau. The Moon probably formed from
debris of a giant impact on the early
Earth.
\item \emph{Mars}: polar caps, huge
volcanoes, channels, alluvial fans and
other signs of past running water; its two
moons are probably captured asteroids.
\item \emph{Jupiter}: the Great Red Spot is
a long-lived storm. Io, a Galilean moon,
is volcanically active.
\item Small moons confine Saturn's rings
and supply them with debris; Titan has a
thick atmosphere.
\item Asteroids lie mostly between Mars
and Jupiter; comets come from the Kuiper
belt and the Oort cloud.
\end{itemize}

A \emph{meteoroid} is a small body in
space; a \emph{meteor} is the streak of
light it makes burning in the atmosphere; a
\emph{meteorite} is what reaches the
ground. Craters are rarer on Earth than on
the Moon because the atmosphere burns up
small bodies and erosion and plate
tectonics erase old craters.

\bibliography{my}
\bibliographystyle{amsalpha}

\end{document}
